\section*{第 \ref{ch:g6:sun-earth-moon} 章 --- 太阳--地球--月球系统}

\begin{solution}{exo:g6:sun-earth-moon:1}
地球绕自己的地轴自转，形成昼夜；地球绕太阳公转，
一周为一年；月球绕地球公转，一周为一个月。
\end{solution}

\begin{solution}{exo:g6:sun-earth-moon:2}
穿过两极、地球绕其自转的假想直线。
它大约倾斜直角的四分之一，并且在整个公转轨道上
始终朝同一方向倾斜——倾斜与
指向不变共同形成了四季。
\end{solution}

\begin{solution}{exo:g6:sun-earth-moon:3}
地球绕日一周需要 $365$ 天又四分之一；积攒四个四分之一天，
便成为闰年多出的一天。没有它，日历会逐渐偏离季节——
大约每四年偏差一天，
一个世纪便偏差数周。
\end{solution}

\begin{solution}{exo:g6:sun-earth-moon:4}
地球到月球的距离约相当于 $30$ 个地球，
太阳表面横跨约相当于 $109$ 个地球。
\end{solution}

\begin{solution}{exo:g6:sun-earth-moon:5}
竹签代替地轴，而真实地轴相对恒星保持倾斜不变。
若让竹签摆动——例如始终朝向灯——模型便显示
无尽的夏天：倾斜不再保持不变，
季节随之消失。
\end{solution}

\begin{solution}{exo:g6:sun-earth-moon:6}
相反的情况是：当我们所在的半球朝向太阳倾斜时，
另一个半球必然背离太阳——一个倾斜无法同时有利于两端。
\end{solution}

\begin{solution}{exo:g6:sun-earth-moon:7}
正对时，手电的光以陡峭而集中的方式落下——
形成一小块明亮而刺目的光斑。倾斜时，同样的光
则涂抹在墙上一大片区域，既稀薄又微弱。
因而朝向光源的倾斜才会得到陡峭而集中的光。
\end{solution}

\begin{solution}{exo:g6:sun-earth-moon:8}
这是月球绕地球的公转：月球运行时，我们从不断变化的角度看它
被太阳照亮的那一半——由蛾眉月到满月再回到蛾眉月。
第三种运动便上演了月相的变化。
\end{solution}

\begin{solution}{exo:g6:sun-earth-moon:9}
第一，北半球的冬季出现时，地球正
略微\emph{较近}于太阳运行——距离的说法方向相反。
第二，两个半球在同一时刻、同一距离下季节相反，
任何相同的距离都做不到这一点；
只有倾斜，使一端有利、另一端不利，
才能做到。
\end{solution}

\begin{solution}{exo:g6:sun-earth-moon:10}
若没有倾斜，公转途中的每个位置都一样：太阳
全年正午都升到同样的高度，各地昼夜始终等长，
一年便失去了四季——只剩一个漫长而永恒的春天，
日晷的弧线凝固不变，
历法中最主要的“戏剧”也就随之取消。
\end{solution}

\begin{solution}{exo:g6:sun-earth-moon:11}
月球绕地球一周用了 $27$ 天，这期间地球也沿着
自己的轨道前进，太阳的方向随之改变，因此月球还需
多走约两天半，才能回到与太阳成一线、形成满月的位置。
“绕一圈”与“成一线”是
两个不同的终点。
\end{solution}

\begin{solution}{exo:g6:sun-earth-moon:12}
极地小镇的日周运动圆很小，由于六月的倾斜，
它可以\emph{完全}位于被照亮的半球内：小镇转动时
从不离开白昼——这就是午夜太阳。十二月时，同样的倾斜
使其小圆完全落入黑暗一侧：整圈旋转却没有日出。
倾斜对各纬度的影响，就是它在两极的绝对影响，
只是更为温和。
\end{solution}

\begin{solution}{pb:g6:sun-earth-moon:1}
\textbf{1.} 每厘米对应 \qty{12800}{km}。
\textbf{2.} 约 \qty{0.25}{cm}——一粒胡椒（或一颗米粒）。
\textbf{3.} 距弹珠 $30 \times 1 = \qty{30}{cm}$。
\textbf{4.} $109$ 厘米——约一米：呼啦圈，或一个大号
沙滩球。
\textbf{5.} $150000000 \div 12800 \approx \qty{11700}{cm}$——
约 \qty{117}{m}。
\textbf{6.} 勉强不行：太阳球须立在
\qty{100}{m} 长场地的远端底线之外约 \qty{17}{m} 处。
\textbf{7.} 约 $117$ 秒——大约两分钟。
\textbf{8.} 太阳：$109 \div 11700 \approx 0.009$；月球小粒：
$0.25 \div 30 \approx 0.008$——几乎相等：二者在弹珠上看起来
大小相同。在我们的天空中也是如此——当月球恰好从太阳前方经过时，
它几乎恰好将太阳遮住：
形成日全食，大自然正好利用了这一巧合。
\textbf{9.} 弹珠绕行一周：约 $365$ 秒——
约六分钟的步行；小粒绕弹珠一周：约
$27$ 秒。
\textbf{10.} 每秒自转一周——在光滑而无标记的
一厘米大的弹珠上，自转看不出任何痕迹，
眼睛无从察觉：在上面画一个点才是切实可行的办法。
\textbf{11.} 从竖直状态倾斜约直角的四分之一——
并在整个绕行过程中始终指向
场地上同一个远处的角落，从不转向
太阳球。
\textbf{12.} 例如：模型真实地表现了距离与大小，
而且\emph{同时}呈现——巨大的空旷正是它的成功之处。
它无法同时展示按真实速度进行的运动（
六分钟的“一年”已经是极度加速；真实的弹珠
将会爬行）——当然，还有光与热：没有
直径一米的球，无法像太阳照亮行星那样
照亮运动场。
\end{solution}
